\documentclass[11pt,a4paper]{article}
% Preâmbulo comum. Não duplicar pacotes nos módulos.
\usepackage[utf8]{inputenc}
\usepackage[T1]{fontenc}
\usepackage[brazil]{babel}
\usepackage{amsmath,amssymb,amsthm,mathtools}
\usepackage{geometry}
\geometry{a4paper,margin=2.3cm}
\usepackage{microtype}
\usepackage{enumitem}
\usepackage{graphicx}
\usepackage{xcolor}
\usepackage{booktabs}
\usepackage{hyperref}
\hypersetup{colorlinks=true,linkcolor=blue!40!black,urlcolor=blue!50!black}

\setcounter{tocdepth}{2}
\numberwithin{equation}{section}

\theoremstyle{plain}
\newtheorem{teorema}{Teorema}[section]
\newtheorem{proposicao}[teorema]{Proposição}
\newtheorem{lema}[teorema]{Lema}
\newtheorem{corolario}[teorema]{Corolário}

\theoremstyle{definition}
\newtheorem{definicao}[teorema]{Definição}
\newtheorem{exemplo}[teorema]{Exemplo}
\newtheorem{exercicio}{Exercício}[section]

\theoremstyle{remark}
\newtheorem{observacao}[teorema]{Observação}
\newtheorem{contraexemplo}[teorema]{Contraexemplo}
\newtheorem{intuicao}[teorema]{Intuição}

\newcommand{\R}{\mathbb{R}}
\newcommand{\N}{\mathbb{N}}
\newcommand{\Z}{\mathbb{Z}}
\newcommand{\Q}{\mathbb{Q}}
\DeclareRobustCommand{\eps}{\varepsilon}
\DeclareRobustCommand{\vet}[1]{\vec{#1}}
\DeclareRobustCommand{\dif}{\mathrm{d}}
\DeclareRobustCommand{\abs}[1]{\left\lvert #1 \right\rvert}
\DeclareRobustCommand{\norm}[1]{\left\lVert #1 \right\rVert}
\DeclareMathOperator{\grad}{grad}
\DeclareMathOperator{\Div}{div}
\DeclareMathOperator{\rot}{rot}

\newcommand{\base}[1]{\par\medskip\noindent\textbf{Base para relembrar.} #1\par\medskip}
\newcommand{\cuidado}[1]{\par\medskip\noindent\textbf{Quando a hipótese falha.} #1\par\medskip}

\newcommand{\figura}[2]{%
  \begin{figure}[htbp]
  \centering
  \IfFileExists{#1}{%
    \includegraphics[width=0.84\textwidth]{#1}%
  }{%
    \fbox{\parbox{0.8\textwidth}{\centering Figura ainda não gerada:\\ \texttt{#1}}}%
  }
  \caption{#2}
  \end{figure}%
}

\newenvironment{esqueleto}{%
  \par\medskip\noindent\textbf{Esqueleto da demonstração.}\par\nopagebreak
  \begin{enumerate}[leftmargin=1.4em]
}{\end{enumerate}\medskip}

\title{Cálculo III: integrais múltiplas}
\author{}
\date{}
\begin{document}
\maketitle
\tableofcontents

\section{Integral dupla no retângulo}

\base{A integral de Riemann em $[a,b]$ é o limite das somas $f(x_i^*)\Delta x_i$ quando a malha da partição tende a zero, para $f$ contínua. O teorema de Fubini reduz a integral dupla a duas integrais de uma variável, numa ordem ou na outra.}

Seja $R=[a,b]\times[c,d]$ e seja $f\colon R\to\R$ limitada. Uma partição de $R$ é um produto de partições de $[a,b]$ e de $[c,d]$: surgem sub-retângulos $R_{ij}$ de área $\Delta A_{ij}=\Delta x_i\Delta y_j$. A \emph{soma de Riemann} associada a pontos amostrais $(x_{ij}^*,y_{ij}^*)\in R_{ij}$ é
\[
\sum_{i,j} f(x_{ij}^*,y_{ij}^*)\,\Delta A_{ij}.
\]
A malha é o maior lado dentre os $\Delta x_i$ e os $\Delta y_j$. Dizemos que $f$ é \emph{integrável} em $R$, com integral $\iint_R f\,\dif A$, quando essas somas convergem para um único número sempre que a malha tende a zero, qualquer que seja a escolha dos pontos amostrais.

\begin{teorema}
Se $f$ é contínua em $R$, então $f$ é integrável em $R$.
\end{teorema}

Não reconstruímos a prova. O fato usado por baixo é a continuidade uniforme de $f$ no compacto $R$: o erro entre somas superiores e inferiores fica menor que qualquer $\eps>0$ quando a malha é fina. Toda conta concreta desta nota parte de uma função contínua numa região fechada e limitada do tipo descrito a seguir, logo integrável.

A integral é linear, e $f\le g$ em $R$ implica $\iint_R f\le\iint_R g$. Se $f\ge 0$, a integral sobre um subconjunto é no máximo a integral sobre o conjunto maior. Essas três propriedades bastam para o confronto da integral gaussiana.

\begin{teorema}[Fubini no retângulo]
Se $f$ é contínua em $R=[a,b]\times[c,d]$, então as duas integrais iteradas existem e coincidem com a integral dupla:
\[
\iint_R f(x,y)\,\dif A
=\int_a^b\Biggl(\int_c^d f(x,y)\,\dif y\Biggr)\dif x
=\int_c^d\Biggl(\int_a^b f(x,y)\,\dif x\Biggr)\dif y.
\]
\end{teorema}

\begin{esqueleto}
\item Para cada $x$ fixo em $[a,b]$, a função $y\mapsto f(x,y)$ é contínua em $[c,d]$, logo a integral interna $g(x)=\int_c^d f(x,y)\,\dif y$ existe.
\item A função $g$ é contínua em $[a,b]$. De fato, pela continuidade uniforme de $f$ em $R$, dado $\eps>0$ existe $\delta>0$ tal que $\abs{x-x_0}<\delta$ implica $\abs{f(x,y)-f(x_0,y)}<\eps/(d-c)$ para todo $y\in[c,d]$, se $d>c$. Então
\[
\abs{g(x)-g(x_0)}\le\int_c^d\abs{f(x,y)-f(x_0,y)}\,\dif y<\eps.
\]
\item Uma função contínua em $[a,b]$ é integrável, então $\int_a^b g(x)\,\dif x$ existe. Comparando somas de Riemann em $x$, depois de integrar em $y$, com somas duplas de $f$, esse número é a integral dupla. A malha fina torna a diferença menor que $\eps$.
\item O mesmo argumento com a ordem trocada produz a outra iterada. As duas coincidem com a integral dupla, logo coincidem entre si.
\end{esqueleto}

O passo que não fingimos ter provado é a identificação completa entre a integral dupla, definida por somas de Riemann bidimensionais, e a iterada. O esqueleto isola o que é continuidade uniforme e o que é passagem ao limite das somas.

\section{Regiões dos tipos I e II}

\begin{definicao}
Uma região $D\subset\R^2$ é do \emph{tipo I} quando existem $a<b$ e funções contínuas $g_1,g_2\colon[a,b]\to\R$, com $g_1\le g_2$, tais que
\[
D=\bigl\{(x,y):a\le x\le b,\ g_1(x)\le y\le g_2(x)\bigr\}.
\]
É do \emph{tipo II} quando
\[
D=\bigl\{(x,y):c\le y\le d,\ h_1(y)\le x\le h_2(y)\bigr\},
\]
com $h_1,h_2$ contínuas e $h_1\le h_2$.
\end{definicao}

Se $f$ é contínua em $D$ e $D$ é dos dois tipos, as integrais iteradas podem ser escritas nas duas ordens. A escolha deixa de ser indiferença quando uma das primitivas internas não é elementar e a outra é.

\begin{exemplo}
A integral
\[
\int_0^1\int_x^1 e^{y^2}\,\dif y\,\dif x
\]
não se calcula na ordem escrita por uma primitiva elementar de $y\mapsto e^{y^2}$. A região é $0\le x\le 1$ e $x\le y\le 1$. Lida na outra ordem, $y$ vai de $0$ a $1$ e, para cada $y$, $x$ vai de $0$ até $y$:
\[
\int_0^1\int_0^y e^{y^2}\,\dif x\,\dif y
=\int_0^1 y\,e^{y^2}\,\dif y.
\]
A substituição $u=y^2$ dá $\dif u=2y\,\dif y$, com $u$ de $0$ a $1$:
\[
\int_0^1 y\,e^{y^2}\,\dif y=\frac12\int_0^1 e^u\,\dif u=\frac12\bigl(e-1\bigr).
\]
O valor da integral dupla é $(e-1)/2$. A troca dos extremos é parte da conta: errar o intervalo de $x$ para cada $y$ produz outra região.
\end{exemplo}

\section{Área, valor médio, massa e centro de massa}

Se $D$ tem área positiva e $f$ é contínua em $D$, a \emph{área} e o \emph{valor médio} são
\[
\mathrm{área}(D)=\iint_D 1\,\dif A,
\qquad
\overline{f}=\frac{1}{\mathrm{área}(D)}\iint_D f\,\dif A.
\]
O valor médio é o número constante que, integrado em $D$, reproduz a integral de $f$.

Uma lâmina ocupando $D$, com densidade contínua $\rho(x,y)\ge 0$ não identicamente nula, tem massa
\[
m=\iint_D \rho(x,y)\,\dif A.
\]
Os momentos em relação aos eixos e o centro de massa $(\bar x,\bar y)$ são
\begin{align*}
M_y&=\iint_D x\,\rho(x,y)\,\dif A,
&
\bar x&=\frac{M_y}{m},\\
M_x&=\iint_D y\,\rho(x,y)\,\dif A,
&
\bar y&=\frac{M_x}{m}.
\end{align*}
O momento $M_y$ leva o fator $x$ porque mede a distribuição na direção horizontal, a que determina $\bar x$. Se a densidade é constante, ela sai da razão e o centro de massa é o centroide geométrico.

\begin{exemplo}
Lâmina no retângulo $[0,1]\times[0,2]$, densidade $\rho(x,y)=x$. A massa é
\[
m=\int_0^2\int_0^1 x\,\dif x\,\dif y=\int_0^2\frac12\,\dif y=1.
\]
O momento que produz $\bar x$ é
\[
M_y=\int_0^2\int_0^1 x^2\,\dif x\,\dif y=\int_0^2\frac13\,\dif y=\frac23,
\]
logo $\bar x=2/3$. O outro momento é
\[
M_x=\int_0^2\int_0^1 yx\,\dif x\,\dif y=\int_0^2 y\cdot\frac12\,\dif y=\frac14\cdot 4=1,
\]
logo $\bar y=1$. A densidade não depende de $y$, então $\bar y$ cai no meio do intervalo $[0,2]$. A densidade cresce com $x$, então $\bar x=2/3$ fica à direita do meio do intervalo $[0,1]$.
\end{exemplo}

\section{Mudança de variáveis}

\begin{teorema}
Sejam $D'\subset\R^2$ uma região dos tipos acima e
\[
x=x(u,v),\qquad y=y(u,v)
\]
funções de classe $C^1$ num aberto que contém $D'$. Suponha a transformação injetora no interior de $D'$ e o jacobiano
\[
J=\det\frac{\partial(x,y)}{\partial(u,v)}
=\det\begin{pmatrix} x_u & x_v \\ y_u & y_v \end{pmatrix}
\]
diferente de zero nesse interior. Se $f$ é contínua em $D$, imagem de $D'$, então
\begin{equation}
\iint_D f(x,y)\,\dif x\,\dif y
=\iint_{D'} f\bigl(x(u,v),y(u,v)\bigr)\,\abs{J(u,v)}\,\dif u\,\dif v.
\label{eq:mudanca}
\end{equation}
\end{teorema}

O módulo é parte da fórmula, não um enfeite. O determinante $J$ tem sinal: ele é positivo quando a transformação preserva a orientação do plano e negativo quando a inverte. A área não enxerga essa orientação. Se alguém integra $J$ sem módulo, uma parametrização que percorre a região ``ao contrário'' produz uma integral negativa de uma função positiva.

\begin{exemplo}
A troca $(u,v)\mapsto(x,y)=(v,u)$ tem
\[
J=\det\begin{pmatrix} 0 & 1 \\ 1 & 0 \end{pmatrix}=-1.
\]
Ela leva o quadrado $[0,1]^2$ nele mesmo e é bijetora. A área do quadrado é $1$. A fórmula com $\abs{J}$ devolve $\iint 1\,\dif u\,\dif v=1$. Sem o módulo, a integral de $J$ seria $-1$, que não é área. O sinal de $J$ registrou que os eixos coordenados foram trocados, invertendo a orientação; a área ignora esse registro.
\end{exemplo}

\section{Coordenadas polares}

A mudança padrão, com $r\ge 0$, é
\[
x=r\cos\theta,\qquad y=r\sin\theta.
\]
O jacobiano é
\begin{align*}
J
&=\det\begin{pmatrix}
\cos\theta & -r\sin\theta \\
\sin\theta & r\cos\theta
\end{pmatrix}
=\cos\theta\cdot r\cos\theta-(-r\sin\theta)\sin\theta
=r(\cos^2\theta+\sin^2\theta)
=r.
\end{align*}
Como $r\ge 0$, tem-se $\abs{J}=r$. No interior do retângulo de parâmetros $r>0$ e $0<\theta<2\pi$, o jacobiano não se anula e a transformação é injetora. A origem, onde $r=0$, e o corte $\theta=0$ identificado com $\theta=2\pi$ têm área zero: não mudam a integral de uma função contínua. Com essa ressalva, a fórmula \eqref{eq:mudanca} vale para discos fechados e fica
\[
\iint_D f(x,y)\,\dif x\,\dif y
=\iint_{D'} f(r\cos\theta,r\sin\theta)\,r\,\dif r\,\dif\theta.
\]
O fator $r$ é a área do retângulo polar deformado: um acréscimo $\Delta r$ e um acréscimo $\Delta\theta$ desenham, longe da origem, um setor de comprimento aproximadamente $r\Delta\theta$ e espessura $\Delta r$.

\figura{../figuras/polar_jacobiano.png}{Retângulo polar e a região correspondente no plano.}

\begin{exemplo}
No disco $x^2+y^2\le 1$,
\begin{align*}
\iint_{x^2+y^2\le 1} e^{-(x^2+y^2)}\,\dif A
&=\int_0^{2\pi}\int_0^1 e^{-r^2}\,r\,\dif r\,\dif\theta.
\end{align*}
A integral em $r$ usa $u=r^2$, $\dif u=2r\,\dif r$:
\[
\int_0^1 e^{-r^2}\,r\,\dif r=\frac12\int_0^1 e^{-u}\,\dif u=\frac12\bigl(1-e^{-1}\bigr).
\]
Multiplicando por $\int_0^{2\pi}\dif\theta=2\pi$,
\[
\iint_{x^2+y^2\le 1} e^{-(x^2+y^2)}\,\dif A=\pi\bigl(1-e^{-1}\bigr).
\]
\end{exemplo}

\section{A integral gaussiana}

\begin{teorema}
\[
\int_{-\infty}^{\infty} e^{-x^2}\,\dif x=\sqrt{\pi}.
\]
\end{teorema}

\begin{proof}
O integrando é positivo, par e contínuo. Defina
\[
I(R)=\int_{-R}^{R} e^{-x^2}\,\dif x,\qquad R>0.
\]
A função $I$ é crescente. A cauda $\int_1^{\infty} e^{-x^2}\,\dif x$ converge por comparação com $\int_1^{\infty} e^{-x}\,\dif x$: para $x\ge 1$ vale $x^2\ge x$, logo $e^{-x^2}\le e^{-x}$. O integrando é par, então a cauda negativa também converge, e a integral imprópria $\int_{-\infty}^{\infty} e^{-x^2}\,\dif x$ existe. O valor é o limite de $I(R)$, porque $I(R)$ é a integral simétrica e as duas caudas somem quando $R\to\infty$. Falta identificar esse limite.

Seja $D_R$ o disco $x^2+y^2\le R^2$ e seja $Q_R=[-R,R]^2$. Vale
\[
D_R\subset Q_R\subset D_{R\sqrt{2}}.
\]
A primeira inclusão: se $x^2+y^2\le R^2$, então $\abs{x}\le R$ e $\abs{y}\le R$. A segunda: se $\abs{x}\le R$ e $\abs{y}\le R$, então $x^2+y^2\le 2R^2$. O integrando $e^{-(x^2+y^2)}$ é positivo, então as integrais respeitam as inclusões:
\[
\iint_{D_R} e^{-(x^2+y^2)}\,\dif A
\le\iint_{Q_R} e^{-(x^2+y^2)}\,\dif A
\le\iint_{D_{R\sqrt{2}}} e^{-(x^2+y^2)}\,\dif A.
\]

Em polares, com o jacobiano $r$ já calculado,
\begin{align*}
\iint_{D_R} e^{-(x^2+y^2)}\,\dif A
&=\int_0^{2\pi}\int_0^R e^{-r^2}\,r\,\dif r\,\dif\theta
=2\pi\cdot\frac12\int_0^{R^2} e^{-u}\,\dif u
=\pi\bigl(1-e^{-R^2}\bigr).
\end{align*}
Trocando $R$ por $R\sqrt{2}$, o disco grande contribui $\pi\bigl(1-e^{-2R^2}\bigr)$.

No quadrado, $f(x,y)=e^{-x^2}e^{-y^2}$ é contínua. Fubini separa o produto:
\[
\iint_{Q_R} e^{-(x^2+y^2)}\,\dif A
=\Biggl(\int_{-R}^{R} e^{-x^2}\,\dif x\Biggr)
\Biggl(\int_{-R}^{R} e^{-y^2}\,\dif y\Biggr)
=I(R)^2.
\]
Juntando as três regiões,
\[
\pi\bigl(1-e^{-R^2}\bigr)
\le I(R)^2
\le\pi\bigl(1-e^{-2R^2}\bigr).
\]
Quando $R\to\infty$, o extremo inferior e o superior tendem a $\pi$. Pelo confronto, $I(R)^2\to\pi$. Como $I(R)>0$, $I(R)\to\sqrt{\pi}$.
\end{proof}

Elevar ao quadrado é o passo que transforma um integral sem primitiva elementar numa integral dupla com simetria rotacional. O confronto entre quadrado e discos é o que autoriza passar ao limite polar: sem ele, integrar em $\R^2$ ``em polares'' seria uma troca de região não justificada.

\section{Integrais triplas}

Se $f$ é contínua na caixa $[a,b]\times[c,d]\times[p,q]$, o teorema de Fubini em três variáveis iguala a integral tripla a qualquer uma das seis ordens iteradas. Não reconstruímos a prova; ela repete a continuidade uniforme, agora com duas integrais internas.

Uma região entre dois gráficos contínuos, sobre uma base $D$ do tipo I ou II no plano $xy$, é
\[
E=\bigl\{(x,y,z):(x,y)\in D,\ u_1(x,y)\le z\le u_2(x,y)\bigr\}.
\]
Para $f$ contínua em $E$,
\[
\iiint_E f\,\dif V
=\iint_D\Biggl(\int_{u_1(x,y)}^{u_2(x,y)} f(x,y,z)\,\dif z\Biggr)\dif A.
\]
O volume é o caso $f=1$.

\begin{exemplo}
O tetraedro $x\ge 0$, $y\ge 0$, $z\ge 0$, $x+y+z\le 1$ tem volume
\begin{align*}
V
&=\int_0^1\int_0^{1-x}\int_0^{1-x-y}\dif z\,\dif y\,\dif x
=\int_0^1\int_0^{1-x}(1-x-y)\,\dif y\,\dif x.
\end{align*}
A integral em $y$ vale
\[
(1-x)y-\frac{y^2}{2}\Big|_{0}^{1-x}=\frac{(1-x)^2}{2}.
\]
Com $u=1-x$,
\[
V=\int_0^1\frac{u^2}{2}\,\dif u=\frac13\cdot\frac12=\frac16.
\]
\end{exemplo}

\section{Cilíndricas e esféricas}

\subsection{Cilíndricas}

A convenção é $r\ge 0$,
\[
x=r\cos\theta,\qquad y=r\sin\theta,\qquad z=z,
\]
e o elemento de volume $r\,\dif r\,\dif\theta\,\dif z$. O jacobiano da transformação $(r,\theta,z)\mapsto(x,y,z)$ é o determinante
\[
\det\begin{pmatrix}
\cos\theta & -r\sin\theta & 0 \\
\sin\theta & r\cos\theta & 0 \\
0 & 0 & 1
\end{pmatrix}.
\]
Expandindo pela terceira linha, a única entrada não nula é o $1$ do canto inferior direito, e o sinal do cofator nessa posição é positivo. O menor correspondente é
\[
\cos\theta\cdot r\cos\theta-(-r\sin\theta)\sin\theta=r.
\]
Assim $J=r$. Com $r\ge 0$, $\abs{J}=r$, e
\[
\iiint_E f(x,y,z)\,\dif V
=\iiint_{E'} f(r\cos\theta,r\sin\theta,z)\,r\,\dif r\,\dif\theta\,\dif z.
\]
Conferência rápida: o cilindro $r\le R$, $0\le z\le H$ tem volume
\[
\int_0^H\int_0^{2\pi}\int_0^R r\,\dif r\,\dif\theta\,\dif z
=H\cdot 2\pi\cdot\frac{R^2}{2}=\pi R^2 H,
\]
que é a fórmula conhecida.

\subsection{Esféricas}

A convenção fica fixa. $\rho\ge 0$ é a distância à origem, $\varphi\in[0,\pi]$ é o ângulo a partir do eixo $z$ positivo, e $\theta\in[0,2\pi)$ é o ângulo no plano $xy$ a partir do eixo $x$. As fórmulas são
\[
x=\rho\sin\varphi\cos\theta,\qquad
y=\rho\sin\varphi\sin\theta,\qquad
z=\rho\cos\varphi.
\]
\figura{../figuras/coordenadas_esfericas.png}{$\rho$, $\varphi$ e $\theta$.}

O jacobiano é o determinante da matriz das derivadas parciais, com colunas na ordem $\partial/\partial\rho$, $\partial/\partial\varphi$, $\partial/\partial\theta$:
\[
\begin{pmatrix}
\sin\varphi\cos\theta & \rho\cos\varphi\cos\theta & -\rho\sin\varphi\sin\theta \\
\sin\varphi\sin\theta & \rho\cos\varphi\sin\theta & \rho\sin\varphi\cos\theta \\
\cos\varphi & -\rho\sin\varphi & 0
\end{pmatrix}.
\]
Expandimos pela terceira linha. O cofator da entrada $\cos\varphi$ é o menor
\begin{align*}
&\det\begin{pmatrix}
\rho\cos\varphi\cos\theta & -\rho\sin\varphi\sin\theta \\
\rho\cos\varphi\sin\theta & \rho\sin\varphi\cos\theta
\end{pmatrix}\\
&\qquad=\rho^2\cos\varphi\sin\varphi\cos^2\theta+\rho^2\sin\varphi\cos\varphi\sin^2\theta
=\rho^2\sin\varphi\cos\varphi.
\end{align*}
A entrada do meio da terceira linha é $-\rho\sin\varphi$, e o cofator correspondente traz um sinal negativo, de modo que a contribuição é
\[
(-\rho\sin\varphi)\cdot(-1)\cdot
\det\begin{pmatrix}
\sin\varphi\cos\theta & -\rho\sin\varphi\sin\theta \\
\sin\varphi\sin\theta & \rho\sin\varphi\cos\theta
\end{pmatrix}.
\]
O menor vale
\begin{align*}
&\sin\varphi\cos\theta\cdot\rho\sin\varphi\cos\theta
-(-\rho\sin\varphi\sin\theta)\sin\varphi\sin\theta
=\rho\sin^2\varphi.
\end{align*}
A contribuição é $\rho\sin\varphi\cdot\rho\sin^2\varphi=\rho^2\sin^3\varphi$. A entrada da direita na terceira linha é zero. Somando,
\begin{align*}
J
&=\cos\varphi\cdot\rho^2\sin\varphi\cos\varphi+\rho^2\sin^3\varphi
=\rho^2\sin\varphi\,(\cos^2\varphi+\sin^2\varphi)
=\rho^2\sin\varphi.
\end{align*}
No intervalo $\varphi\in[0,\pi]$ tem-se $\sin\varphi\ge 0$, e $\rho\ge 0$, logo
\[
\abs{J}=\rho^2\sin\varphi.
\]
O elemento de volume é $\rho^2\sin\varphi\,\dif\rho\,\dif\varphi\,\dif\theta$. Fora desse intervalo de $\varphi$ o seno troca de sinal e o módulo deixa de ser opcional; dentro da convenção fixa, $\rho^2\sin\varphi$ já é não negativo.

\section{Volume da bola e centro de massa do hemisfério}

\begin{exemplo}
A bola $x^2+y^2+z^2\le R^2$ é descrita por $0\le\rho\le R$, $0\le\varphi\le\pi$ e $0\le\theta<2\pi$. O volume é
\begin{align*}
V
&=\int_0^{2\pi}\int_0^{\pi}\int_0^R \rho^2\sin\varphi\,\dif\rho\,\dif\varphi\,\dif\theta
=\Biggl(\int_0^{2\pi}\dif\theta\Biggr)
\Biggl(\int_0^{\pi}\sin\varphi\,\dif\varphi\Biggr)
\Biggl(\int_0^R\rho^2\,\dif\rho\Biggr)\\
&=2\pi\cdot\bigl[-\cos\varphi\bigr]_0^{\pi}\cdot\frac{R^3}{3}
=2\pi\cdot\bigl(1-(-1)\bigr)\cdot\frac{R^3}{3}
=\frac{4}{3}\pi R^3.
\end{align*}
A integral em $\varphi$ vale $2$ porque $\varphi$ percorre o intervalo inteiro $[0,\pi]$, do polo norte ao polo sul.
\end{exemplo}

\begin{exemplo}
Hemisfério sólido superior, $x^2+y^2+z^2\le R^2$ e $z\ge 0$, de densidade constante. A constante cancela na razão do centro de massa; fazemos a conta com densidade $1$. Por simetria em relação aos planos $x=0$ e $y=0$, $\bar x=\bar y=0$: o integrando $x$ é ímpar em $x$ e a região é simétrica, e o mesmo vale para $y$.

O volume é metade do volume da bola,
\[
m=\frac12\cdot\frac{4}{3}\pi R^3=\frac{2}{3}\pi R^3.
\]
Em esféricas, $z=\rho\cos\varphi$ e o hemisfério superior é $0\le\varphi\le\pi/2$. O momento é
\begin{align*}
\iiint z\,\dif V
&=\int_0^{2\pi}\int_0^{\pi/2}\int_0^R (\rho\cos\varphi)\,\rho^2\sin\varphi\,\dif\rho\,\dif\varphi\,\dif\theta\\
&=2\pi\cdot\frac{R^4}{4}\cdot\int_0^{\pi/2}\sin\varphi\cos\varphi\,\dif\varphi.
\end{align*}
Com $u=\sin\varphi$, $\dif u=\cos\varphi\,\dif\varphi$, $u$ vai de $0$ a $1$ e a integral vale $1/2$. Logo
\[
\iiint z\,\dif V=2\pi\cdot\frac{R^4}{4}\cdot\frac12=\frac{\pi R^4}{4}.
\]
A cota do centro de massa é
\[
\bar z=\frac{\pi R^4/4}{2\pi R^3/3}=\frac{\pi R^4}{4}\cdot\frac{3}{2\pi R^3}=\frac{3R}{8}.
\]
\end{exemplo}

\section{Exercícios}

\begin{exercicio}[{[Básico --- iterada no triângulo]}]
Calcule $\displaystyle\iint_D (x+2y)\,\dif A$, em que $D$ é o triângulo $0\le x\le 1$, $0\le y\le 1-x$.
\end{exercicio}

\begin{exercicio}[{[Básico --- Fubini nas duas ordens]}]
Calcule $\displaystyle\int_0^1\int_0^2 xy\,\dif y\,\dif x$ e repita a conta na ordem $\dif x\,\dif y$. Diga por que o teorema de Fubini se aplica.
\end{exercicio}

\begin{exercicio}[{[Básico --- polar direta]}]
Calcule $\displaystyle\iint_{x^2+y^2\le 1}(x^2+y^2)\,\dif A$ em polares, com o fator $r$.
\end{exercicio}

\begin{exercicio}[{[Básico --- valor médio]}]
Calcule o valor médio de $f(x,y)=x^2+y^2$ no disco unitário $x^2+y^2\le 1$.
\end{exercicio}

\begin{exercicio}[{[Intermediário --- trocar a ordem]}]
Calcule $\displaystyle\int_0^1\int_y^1 e^{x^2}\,\dif x\,\dif y$ trocando a ordem. A primitiva interna de $e^{x^2}$ em $x$ não é elementar; a outra ordem é.
\end{exercicio}

\begin{exercicio}[{[Intermediário --- centro de massa]}]
Lâmina homogênea no quarto de disco $x\ge 0$, $y\ge 0$, $x^2+y^2\le 1$. Calcule a massa e o centro de massa. A densidade vale $1$.
\end{exercicio}

\begin{exercicio}[{[Intermediário --- cilíndricas]}]
Calcule o volume do sólido $x^2+y^2\le 1$, $0\le z\le 1-x^2-y^2$, em coordenadas cilíndricas.
\end{exercicio}

\begin{exercicio}[{[Avançado --- setor esférico]}]
Calcule o volume da região $\rho\le R$ e $0\le\varphi\le\pi/4$, com $\theta\in[0,2\pi)$. Na convenção desta nota, $\varphi=\pi/4$ é o cone $z=\sqrt{x^2+y^2}$. Use o jacobiano $\rho^2\sin\varphi$ e simplifique o resultado.
\end{exercicio}

\begin{exercicio}[{[Avançado --- centro de massa do cone]}]
Cone sólido de vértice na origem, eixo no eixo $z$ e base de raio $R$ no plano $z=H>0$: isto é, $0\le z\le H$ e $x^2+y^2\le (Rz/H)^2$. Densidade constante. Mostre que $\bar x=\bar y=0$ e $\bar z=3H/4$.
\end{exercicio}

\end{document}
